% FILE: 04_implementation_surface.tex
% PROJECT: phd-thesis
% MANUFACTURED: 2026-06-01
% AUTHORITY: ledgers/PROJECT_INDEX.yaml, ledgers/WORKFLOW_RECEIPT.yaml,
%            phd-thesis/projects/phd-thesis/project_manifest.yaml
% DO NOT CLAIM: See ledgers/DO_NOT_CLAIM_LEDGER.md

\section{Implementation Surface}
\label{sec:phd-thesis-impl}

\subsection{Source Files}
\label{sec:phd-thesis-impl-sources}

The populated source files of the \texttt{phd-thesis} workspace at the time of this writing
are confined to the \texttt{ledgers/} and \texttt{projects/} directories. The following
files are confirmed present and were read as primary evidence for this corpus:

\begin{itemize}
  \item \textbf{\texttt{phd-thesis/README.md}} --- Workspace purpose, directory structure,
    ALIVE criteria, and doctrine constraints. The authoritative declaration of what this
    workspace is and what conditions govern its release.

  \item \textbf{\texttt{ledgers/WORKFLOW\_RECEIPT.yaml}} --- Manufacturing provenance record.
    Binds the workspace to branch \texttt{phd-thesis-corpus-manufacture-001}, pre-manufacture
    commit \texttt{b4279892881e5a8ff83f31b273219a0740fead0b}, and a verified toolchain
    (pdfTeX 3.141592653, Python 3.14.3, git 2.51.2, ripgrep 15.1.0). Phase:
    \texttt{SETUP\_COMPLETE}.

  \item \textbf{\texttt{ledgers/DO\_NOT\_CLAIM\_LEDGER.md}} --- Claim prohibition registry.
    Seven absolute prohibitions governing all artifacts manufactured in this workspace.

  \item \textbf{\texttt{ledgers/PROJECT\_INDEX.yaml}} --- Canonical ledger of 34 research
    projects. Each entry records: \texttt{slug}, \texttt{absolute\_path},
    \texttt{manifest} path, \texttt{description}, \texttt{detected\_languages},
    \texttt{detected\_frameworks}, \texttt{detected\_research\_surfaces},
    \texttt{likely\_thesis\_role}, \texttt{key\_files}, and boolean flags
    \texttt{readme\_present}, \texttt{receipt\_present}, \texttt{ontology\_present}.

  \item \textbf{\texttt{projects/phd-thesis/project\_manifest.yaml}} --- The workspace's
    own manifest, declaring its role as \texttt{dissertation-corpus-manufacturing-workspace}.

  \item \textbf{\texttt{projects/ggen/project\_manifest.yaml}} --- Manifest for the
    Governance Generation Engine, which is the intended chapter manufacturing engine.

  \item \textbf{\texttt{projects/checkpoints/project\_manifest.yaml}} --- Manifest for the
    checkpoints corpus (sealed verdicts: PROCESS\_INTELLIGENCE\_ALIVE\_001,
    PI\_RESEARCH\_PROGRAM\_ALIVE\_001, GGEN\_ECOSYSTEM\_INTEL\_ALIVE\_001,
    SUBSTRATE\_COMPLETE\_001).

  \item \textbf{\texttt{projects/receipts/project\_manifest.yaml}} --- Manifest for the
    receipt chain corpus (5 JSON receipts with ed25519 signatures and BLAKE3-intended chains).

  \item \textbf{\texttt{projects/research-pi-program/project\_manifest.yaml}} --- Manifest
    for the PI Program governance layer (19 TTL files, 60+ SPARQL queries, manufacturing
    receipts, autonomic activation logs).
\end{itemize}

The source files that are \emph{not yet present} but declared in the directory structure are:
all files under \texttt{chapters/}, \texttt{frontmatter/}, \texttt{build/}, and
\texttt{scripts/}. These are empty stubs awaiting chapter manufacture.

\subsection{Commands}
\label{sec:phd-thesis-impl-commands}

No manufacturing scripts are present in the \texttt{scripts/} directory at the time of this
writing (phase: \texttt{SETUP\_COMPLETE}). The toolchain for the next phase has been verified
and recorded in \texttt{ledgers/WORKFLOW\_RECEIPT.yaml}:

\begin{itemize}
  \item \textbf{PDF compilation}: \texttt{pdflatex} (pdfTeX 3.141592653-2.6-1.40.26,
    TeX Live 2024). The declared command surface is \texttt{pdflatex master.tex} from the
    \texttt{build/} directory, though the master file has not yet been manufactured.
  \item \textbf{Receipt hashing}: The receipt doctrine requires BLAKE3 hashing of the
    compiled PDF. The BLAKE3 toolchain availability is not yet recorded in the workflow
    receipt; this is a gap for the next phase.
  \item \textbf{Ontology queries}: \texttt{ggen} (the Governance Generation Engine) is the
    declared manufacturing framework. Its invocation pattern --- SPARQL queries over Turtle
    ontologies, rendered through Tera templates --- is documented in the ggen project manifest
    and in \texttt{ggen/ggen.toml}.
  \item \textbf{Version control}: \texttt{git 2.51.2}. All manufactured artifacts are
    committed on branch \texttt{phd-thesis-corpus-manufacture-001} under conventional commit
    format with research-specific scopes as defined in the process-intelligence CLAUDE.md.
\end{itemize}

The clap-noun-verb command grammar referenced in the thesis lineage
(\texttt{ledgers/PROJECT\_INDEX.yaml} lineage header) is the CLI design pattern used in the
wasm4pm and ggen command surfaces. It is not directly instantiated in the \texttt{phd-thesis}
workspace but governs the command interface of the manufacturing tools consumed by this
workspace.

\subsection{Data and Ontology Surfaces}
\label{sec:phd-thesis-impl-data}

The data surface of this workspace is the \texttt{ledgers/PROJECT\_INDEX.yaml}. Its schema
is YAML with a \texttt{projects:} array; each entry is a structured record as described in
Section~\ref{sec:phd-thesis-impl-sources}. The index is the primary data artifact of the
workspace: 34 entries, manufactured 2026-06-01.

The ontology surfaces consumed by this workspace --- but not hosted within it --- are
distributed across the foundry:

\begin{itemize}
  \item \textbf{PI Program ontology}: 19 TTL files in
    \texttt{research/pi-program/ggen/ontology/}, covering pi-program.ttl,
    governance-policy.ttl, and autonomic-law.ttl. These are the SPARQL-queryable authority
    sources for the manufacturing pipeline.
  \item \textbf{Prompt manufactory ontology}: TTL files in
    \texttt{research/prompt-manufactory/ggen/ontology/}, covering prompt-manufactory.ttl,
    workflow-law.ttl, and skill-law.ttl.
  \item \textbf{ggen ontology extensions}: \texttt{ggen/ontology-extensions.ttl} and
    \texttt{ggen/wasm4pm-compat.ttl}.
  \item \textbf{Open ontologies}: PROV-O, SKOS, ODRL, DCTERMS, SHACL mappings in
    \texttt{research/open-ontologies/federated/CROSS\_PROJECT\_DEPENDENCY\_GRAPH\_20260601.yaml}.
\end{itemize}

The workspace does not host an ontology directly (\texttt{ontology\_present: false} in the
project manifest). It consumes ontology surfaces through the manufacturing tools.

\subsection{Templates and Rendered Artifacts}
\label{sec:phd-thesis-impl-templates}

No Tera templates are hosted within the \texttt{phd-thesis} workspace. The template surface
is distributed across the foundry and will be consumed during chapter manufacture:

\begin{itemize}
  \item \textbf{\texttt{ggen/templates/blue-river.tera}} --- Template for the Blue River Dam
    orchestrator artifact.
  \item \textbf{\texttt{templates/pcp-types.ts.tera}} --- Template for Post-Cyberpunk PCP
    TypeScript type surface.
  \item \textbf{Prompt manufactory templates} in
    \texttt{research/prompt-manufactory/ggen/templates/} --- Templates for checkpoint,
    hook, skill, subagent, and workflow prompt artifacts.
\end{itemize}

The TeX files manufactured in this corpus (including this file) are the first rendered
artifacts of the \texttt{phd-thesis} workspace's chapter manufacturing phase. They are
manufactured directly rather than through a ggen template pipeline, because no ggen.toml
for the thesis chapter pipeline has been confirmed present in the workspace as of
2026-06-01 (this is an open question documented in the gaps). The ggen.toml that governs
the manufacturing infrastructure itself exists at \texttt{ggen/ggen.toml}, but no
thesis-chapter-specific configuration has been recorded.
